\Hefteintrag{1}{Datenschutz}{
Datenschutz beschreibt den \emphColA{Schutz vor der missbräuchlichen Verarbeitung 
personenbezogener Daten} sowie den Schutz des Rechts auf informationelle 
Selbstbestimmung, also dass jeder jederzeit darüber bestimmt, wer wann seine Daten nutzen und verarbeiten kann.
}

\Hefteintrag{1}{Personenbezogene Daten}{
 Personenbezogene Daten sind \emphColA{alle Informationen, die sich auf eine} (identifizierbare) 
\emphColA{Person beziehen}. 
Personenbezogene Daten, die in einer Weise anonymisiert worden sind, dass die 
betroffene Person nicht oder nicht mehr identifiziert werden kann, gelten nicht mehr 
als personenbezogene Daten. Damit die Daten \emphColB{wirklich anonymisiert} sind, muss die 
\emphColB{Anonymisierung unumkehrbar} sein.

\vspace{0.5cm}

Dem gegenüber steht die \emphColC{Pseudonymisierung}, die z.B. mit dem Ergänzen einer weiteren Tabellenspalte \emphColC{rückgängig gemacht werden kann}.
}